\documentclass[11pt]{amsart}
\usepackage[margin=1.1in]{geometry}
\usepackage{lmodern}
\usepackage{amsmath,amssymb,amsthm,mathtools}
\usepackage[mathscr]{eucal}
\usepackage{booktabs,enumitem,xcolor,array}
\usepackage[colorlinks=true,linkcolor=blue!50!black,citecolor=green!40!black,urlcolor=blue!60!black]{hyperref}
\hypersetup{
  pdftitle={Hermite interpolation at the prime powers: a positivity calculus, a polynomial model, and a family of test functions vanishing at the zeros of the Riemann zeta function},
  pdfauthor={Nic Johns},
  pdfkeywords={Riemann zeta function, Weil explicit formula, Hermite interpolation, Christoffel identity, Fourier interpolation, positivity, computer-assisted proof}
}

\setlength{\emergencystretch}{2.5em}
% macros.tex -- shared macros for main.tex and sections/*.tex.
% The notation follows the companion paper (positive solutions of the explicit formula).
% New local macros in a section file: use \providecommand and a section-specific prefix.

% ---------------------------------------------------------------- general
\newcommand{\RR}{\mathbb R}
\newcommand{\ZZ}{\mathbb Z}
\newcommand{\CC}{\mathbb C}
\newcommand{\NN}{\mathbb N}
\newcommand{\QQ}{\mathbb Q}
\newcommand{\dd}{\,\mathrm d}
\newcommand{\eps}{\varepsilon}
\renewcommand{\Re}{\operatorname{Re}}
\renewcommand{\Im}{\operatorname{Im}}
\newcommand{\supp}{\operatorname{supp}}
\newcommand{\sgn}{\operatorname{sgn}}
\newcommand{\Wr}{\operatorname{Wr}}
\newcommand{\one}{\mathbf 1}
\newcommand{\FT}[1]{\widehat{#1}}          % Fourier transform, \int f(t)e^{-2\pi i\xi t}dt
\newcommand{\hF}{\widehat F}
\newcommand{\hyp}[1]{\textup{(#1)}}       % named hypothesis tag, e.g. \hyp{N}, \hyp{Mg}

% ---------------------------------------------------------------- coordinates
\newcommand{\xin}[1]{\xi_{#1}}            % \xi_n = \log n/(2\pi)
\newcommand{\Eul}{\theta}                 % Euler operator x d/dx = y d/dy

% ---------------------------------------------------------------- arithmetic data
\newcommand{\PP}{\mathrm{PP}}             % prime powers > 1; \PP_J the first J
\newcommand{\xiPP}{\xi_{\mathrm{PP}}}     % {\xi_n : n \in PP}
\newcommand{\BRS}{\mathcal M}             % {\log m/(4\pi) : m \ge 1}
\newcommand{\GR}{\Gamma_{\mathbb R}}
\newcommand{\GC}{\Gamma_{\mathbb C}}
\newcommand{\Xiinf}{\Xi_\infty}           % archimedean part of \Xi
\newcommand{\Zz}{Z_\zeta}                 % real zero ordinates of zeta
\newcommand{\muz}{\mu_\zeta}
\newcommand{\nuz}{\nu_\zeta}
\newcommand{\pz}{p_\zeta}

% ---------------------------------------------------------------- framework
\newcommand{\Oinf}{\Omega_\infty}         % archimedean density
\newcommand{\Arch}{\mathcal A}            % archimedean functional; \Arch_q with conductor q
\newcommand{\Tc}{\mathcal T}              % test class; \Tc_\delta
\newcommand{\Gc}{\mathcal G}              % Gaussian wave packets
\newcommand{\Kad}{\mathcal K}             % admissible pairs
\newcommand{\Cone}{\mathcal C}
\newcommand{\ConeOPS}{\mathcal C_{\mathrm{OPS}}}
\newcommand{\kstar}{\kappa^{*}}
\newcommand{\kOPS}{\kappa^{*}_{\mathrm{OPS}}}
\newcommand{\qmin}{q_{\min}}
\newcommand{\condE}{\textup{(E)}}
\newcommand{\condU}{\textup{(U)}}
\newcommand{\condS}{\textup{(S)}}
\newcommand{\RH}{\mathrm{RH}}

% ------------------------------------------------------------- zero-killing family
\newcommand{\Wnull}{W_0}                  % z^2[(z^2+9)K_0 - 6zK_1]
\newcommand{\Psione}{\Psi_1}              % Gamma-only kernel
\newcommand{\kt}{\tilde\kappa}            % sqrt(2z/pi) e^z K_0(z)
\newcommand{\Hc}{\mathscr H}              % growth class \Hc_a
\newcommand{\Jset}{\mathcal J}            % {J : M_J non-singular}
\newcommand{\Frep}{F_{\mathrm{rep}}}      % certified rational representative
\newcommand{\TEL}{\mathrm{TEL}}
\newcommand{\PT}{\mathrm{PT}}

% ---------------------------------------------------------------- interpolation calculus
\newcommand{\Fc}{\mathcal F}              % \Fc_P = P(-\theta^2)\varphi
\newcommand{\Asp}{\mathscr A}             % nested spaces \Asp_K
\newcommand{\yv}{\mathbf y}               % node set (y_1<...<y_J)
\newcommand{\yPP}{\mathbf y^{\mathrm{PP}}}% prime-power node set, \yPP_J

% ---------------------------------------------------------------- the polynomial model
\newcommand{\Hank}{\mathcal H_0}            % Hankel-type transform
\newcommand{\Lag}{\mathfrak L}               % y^m -> m! L_m(y)
\newcommand{\Lrule}{\mathcal L}              % the p_1-rule
\newcommand{\Tu}{\operatorname{Tu}}          % Turan ratio
\newcommand{\Rp}{\textup{(R$^{+}$)}}
\newcommand{\FRt}{\textup{(FR)}}
\newcommand{\yle}[1]{\mathbf y_{\le #1}}
\newcommand{\pf}{\mathsf f}                  % monic cofactor of a chain
\newcommand{\pg}{\mathsf g}                  % 0-augmented cofactor


% ---------------------------------------------------------------- environments
% theorem/proposition/lemma/corollary: proved here (possibly computer-assisted, with the
% computation and ancillary file named).  conjecture: open.  hypothesis: a named open
% condition used as an assumption.  observation ("Numerical observation"): computed, not
% certified, never used in a proof; all of them are in Section~\ref{sec:observations}.
\newtheorem{theorem}{Theorem}[section]
\newtheorem{proposition}[theorem]{Proposition}
\newtheorem{lemma}[theorem]{Lemma}
\newtheorem{corollary}[theorem]{Corollary}
\newtheorem{conjecture}[theorem]{Conjecture}
\newtheorem{hypothesis}[theorem]{Hypothesis}
\newtheorem{question}[theorem]{Question}
% Statements quoted from the companion paper [I] (no proof here).
\newtheorem{fact}[theorem]{Fact}
\theoremstyle{definition}
\newtheorem{definition}[theorem]{Definition}
\theoremstyle{remark}
\newtheorem{remark}[theorem]{Remark}
\newtheorem{example}[theorem]{Example}
\newtheorem{observation}[theorem]{Numerical observation}
\numberwithin{equation}{section}
% Restatements of results of the body in the introduction: \begin{restate}{Theorem~\ref{...}} ... \end{restate}.
\theoremstyle{plain}
\newcommand{\restatename}{}
\newtheorem*{restateinner}{\restatename}
\newenvironment{restate}[1]{\renewcommand{\restatename}{#1}\begin{restateinner}}{\end{restateinner}}

\newcommand{\inputsection}[1]{\input{sections/#1}}

\title[Hermite interpolation at the prime powers]{Hermite interpolation at the prime powers:\\
a positivity calculus, a polynomial model, and a family\\ of test functions vanishing at the zeros of $\zeta(s)$}

\author{Nic Johns}

\date{October 2026}

\subjclass[2020]{Primary 11M26; Secondary 41A05, 42A82, 33C10, 65G20}
\keywords{Riemann zeta function, Weil explicit formula, Hermite interpolation, Christoffel identity, Fourier
interpolation, positivity, computer-assisted proof}

\begin{document}

\begin{abstract}
We study nested Hermite interpolation with double nodes at the points $y_j=2\pi n_j$, where $n_j$ runs through the
prime powers, for three kernels: the Fourier transform of $\Xi^2$, which defines an explicit family of test functions
$\Xi(t)^2P_J(t^2)$ for the explicit formula of the Riemann zeta function, all vanishing at every non-trivial zero; its archimedean part; and the
kernel $e^{-y}$, which gives a polynomial model, equivalently a radial Fourier interpolation problem on the circles
$\lvert v\rvert=\sqrt{2n}$ in $\RR^2$. A kernel-free calculus (Christoffel, ladder, insertion and half-step identities)
reduces positivity of the interpolant beyond the nodes, in the model, to one likelihood-ratio law per level. A
computer-assisted certificate verifies this law at every level up to $160$, which proves the positivity of the front of
the model for every $J\le80$. We prove far-node asymptotics for the model and for the true kernel, show that sign
regularity of the associated five-function systems cannot give the tail of the model, and show that a natural sign
property fails at the critical node density. For the true family we reduce the open hypotheses of a criterion for
uniqueness and sharpness, proved in a companion paper, to explicit step laws, which remain open. Numerical observations
are collected in one section and are never used in proofs.

\end{abstract}

\maketitle

\setcounter{tocdepth}{1}
\tableofcontents

\inputsection{01-intro}
\inputsection{02-setting}
\inputsection{03-nested}
\inputsection{04-chain}
\inputsection{05-model}
\inputsection{06-model-chain}
\inputsection{07-model-tail}
\inputsection{08-signW}
\inputsection{09-true}
\inputsection{10-observations}
\inputsection{11-conjectures}

\section*{AI-use statement}
This manuscript was produced with substantial use of AI. The mathematical arguments, the numerical computations and the text were generated by instances of Claude Opus 5.5
(Anthropic) working as coordinated agents: one set of agents developed the arguments, separate agents acted as
independent hostile referees of each component, and the referees' requested repairs were incorporated. The zero-side support
theorem of the companion paper, which this manuscript cites, was contributed by Astra (OpenAI) and checked by independent referee
agents before inclusion. The numerical
certificates were computed by AI-written scripts that are included as supplementary material and can be rerun. The
author initiated and directed this project, designed its verification process, and reviewed the results and the
verification record, but did not independently check every step of the written proofs. The author takes full
responsibility for the content of this manuscript, including any errors, and welcomes independent verification and
corrections.

\appendix
\inputsection{A-proofs}
\inputsection{B-certificates}
\inputsection{C-anc}

\bibliographystyle{amsplain-doi}
\bibliography{refs}

\end{document}
